\documentclass[12pt,a4paper]{article}

% imscribing and fonts
\usepackage[utf8]{inputenc}
\usepackage[T1]{fontenc}
\usepackage{lmodern}
\usepackage{textcomp}

% Page layout
\usepackage[top=1in, bottom=1in, left=1in, right=1in]{geometry}

% Typography
\usepackage{microtype}
\usepackage{parskip}

% Language
\usepackage[english]{babel}

% Headers and footers
\usepackage{fancyhdr}
\pagestyle{fancy}
\fancyhf{}
\fancyhead[L]{\leftmark}
\fancyhead[R]{\thepage}
\fancyfoot[C]{\thepage}
\renewcommand{\headrulewidth}{0.4pt}
\renewcommand{\footrulewidth}{0pt}

% Section styling
\usepackage{titlesec}
\titleformat{\section}{\Large\bfseries}{\thesection}{1em}{}
\titleformat{\subsection}{\large\bfseries}{\thesubsection}{1em}{}
\titleformat{\subsubsection}{\normalsize\bfseries}{\thesubsubsection}{1em}{}
\titlespacing*{\section}{0pt}{2.5ex plus 1ex minus .2ex}{1.5ex plus .2ex}
\titlespacing*{\subsection}{0pt}{2ex plus 1ex minus .2ex}{1ex plus .2ex}
\titlespacing*{\subsubsection}{0pt}{1.5ex plus 1ex minus .2ex}{0.8ex plus .2ex}

% List formatting
\usepackage{enumitem}
\setlist[itemize]{topsep=0.5ex, itemsep=0.25ex, parsep=0pt}
\setlist[enumerate]{topsep=0.5ex, itemsep=0.25ex, parsep=0pt}

% Essential packages
\usepackage{graphicx}
\usepackage[table]{xcolor}
\usepackage{booktabs}
\usepackage{array}
\usepackage{tabularx}
\usepackage{longtable}
\usepackage{amsmath}
\usepackage{amssymb}
\usepackage[normalem]{ulem}
\rowcolors{2}{gray!10}{white}
\renewcommand{\arraystretch}{1.3}

% Cross-references
\usepackage{hyperref}
\usepackage{cleveref}

% Custom commands
\newcommand{\pilcrow}{\S}

% Hyperref setup
\hypersetup{
    colorlinks=true,
    linkcolor=blue,
    filecolor=magenta,
    urlcolor=cyan,
}

% Code listings
\usepackage{listings}
\definecolor{codebg}{RGB}{248,248,248}
\definecolor{codeframe}{RGB}{200,200,200}
\definecolor{codekeyword}{RGB}{0,0,180}
\definecolor{codecomment}{RGB}{0,128,0}
\definecolor{codestring}{RGB}{163,21,21}
\lstset{
    basicstyle=\ttfamily\small,
    breaklines=true,
    breakatwhitespace=false,
    frame=single,
    framerule=0.5pt,
    rulecolor=\color{codeframe},
    backgroundcolor=\color{codebg},
    keepspaces=true,
    numbers=left,
    numberstyle=\tiny\color{gray},
    numbersep=8pt,
    showstringspaces=false,
    tabsize=4,
    xleftmargin=1em,
    xrightmargin=1em,
    aboveskip=1em,
    belowskip=1em,
    keywordstyle=\color{codekeyword}\bfseries,
    commentstyle=\color{codecomment}\itshape,
    stringstyle=\color{codestring},
    literate={_}{{\textunderscore}}1
             {-}{{\textminus}}1
             {<}{{\textless}}1
             {>}{{\textgreater}}1
             {~}{{\textasciitilde}}1
             {^}{{\textasciicircum}}1
}

% YAML language definition for listings
\lstdefinelanguage{YAML}{
    keywords={true,false,null,y,n},
    keywordstyle=\color{blue}\bfseries,
    sensitive=false,
    comment=[l]{\#},
    morecomment=[s]{/*}{*/},
    commentstyle=\color{gray}\ttfamily,
    stringstyle=\color{red}\ttfamily,
    morestring=[b]'',
    morestring=[b]'',
}

% TOML language definition for listings
\lstdefinelanguage{TOML}{
    keywords={true,false},
    keywordstyle=\color{blue}\bfseries,
    sensitive=true,
    comment=[l]{\#},
    commentstyle=\color{gray}\ttfamily,
    stringstyle=\color{red}\ttfamily,
    morestring=[b]'',
    morestring=[b]'',
}

% Dockerfile language definition for listings
\lstdefinelanguage{Dockerfile}{
    keywords={FROM,RUN,CMD,LABEL,EXPOSE,ENV,ADD,COPY,ENTRYPOINT,VOLUME,USER,WORKDIR,ARG,ONBUILD,STOPSIGNAL,HEALTHCHECK,SHELL},
    keywordstyle=\color{blue}\bfseries,
    sensitive=false,
    comment=[l]{\#},
    commentstyle=\color{gray}\ttfamily,
    stringstyle=\color{red}\ttfamily,
    morestring=[b]'',
    morestring=[b]'',
}

% JSON language definition for listings
\lstdefinelanguage{JSON}{
    keywords={true,false,null},
    keywordstyle=\color{blue}\bfseries,
    sensitive=true,
    commentstyle=\color{gray}\ttfamily,
    stringstyle=\color{red}\ttfamily,
    morestring=[b]'',
}

% Code listing float environment
\usepackage{float}
\newfloat{listing}{htbp}{lol}[section]
\floatname{listing}{Listing}

% Admonition boxes
\usepackage{tcolorbox}
\tcbuselibrary{skins,breakable}

% Admonition style definitions
\newtcolorbox{admonitionbox}[2][]{
    enhanced,
    breakable,
    colback=#2!5,
    colframe=#2!80!black,
    fonttitle=\bfseries,
    title=#1,
    left=2mm,
    right=2mm,
}

% Imscription tuple boxes
\newtcolorbox{imscriptionbox}{
    enhanced,
    colback=blue!5,
    colframe=blue!75!black,
    fontupper=\large,
    center upper,
    boxrule=1pt,
    arc=4pt,
}

\begin{document}

\title{The Aether and Its Vessel - Two Proofs That $E_8$ Finds Its Perfect Vessel in $G_2$}
\date{\today}

\maketitle

\subsection{A Dual Proof --- Conventional Lie Theory and the Imscribing Grammar}

\begin{center}
\rule{0.5\textwidth}{0.4pt}
\end{center}

\subsection{Abstract}

We present a dual proof that the exceptional Lie group $E_8$ --- the \emph{Aether}, the maximal exceptional structure --- necessarily finds its completion in $G_2$, the smallest exceptional Lie group, which we term the \emph{Vessel}. The proof is given first in the conventional language of Lie theory, root system analysis, and octonion algebra. It is then translated into the Imscribing Grammar (IG), a structural formalism that imscribes physical and mathematical systems as 12-primitive tuples drawn from the crystal of $4 \times 5 \times 4 \times 5 \times 3 \times 5 \times 3 \times 4 \times 5 \times 4 \times 3 \times 4 = 17{,}280{,}000$ structural types. Within the IG, we compute the distance, meet, join, tensor product, ouroboricity tier, consciousness score, promotion signature, principal decomposition, and crystal address for both $E_8$ and $G_2$, demonstrating that the relationship is not merely categorical but \emph{structurally exact}: the Aether--Vessel duality is imscribed with mathematical precision across all 12 primitives.

The IG imscribing reveals that $G_2$ and $E_8$ share 5 primitives at lockstep ($T_{\text{bullseye}}$, $R_{\text{lyoghlig}}$, $F_{\text{hardsign}}$, $K_{\text{schwa}}$, $\Phi_{\text{ctyogh}}$), diverge across 7 with a weighted Euclidean distance of 4.12, and that the promotion path from Vessel to Aether requires 6 structural lifts (including $D_{\text{turnthree}} \to D_{\text{invomega}}$, $\Gamma_{\text{corner}} \to \Gamma_{\text{secstress}}$, $H_0 \to H_2$, $\Omega_{\text{closeepsilon}} \to \Omega_{\text{dzlig}}$). The meet --- the shared structural floor --- captures precisely what the Vessel contributes as the irreducible core. The join --- the minimal ceiling containing both --- reconstructs the Aether's full tuple plus the Vessel's $P_{\text{pipevar}}$ symmetry, revealing the structural surplus that $G_2$ contributes \emph{beyond} what $E_8$ alone demands.

\begin{center}
\rule{0.5\textwidth}{0.4pt}
\end{center}

\subsection{Part I: The Conventional Proof}

\subsubsection{Introduction: The Exceptional Chain}

The exceptional Lie groups form a chain of increasing dimension and complexity:

\[G_2 \subset F_4 \subset E_6 \subset E_7 \subset E_8\]

This chain exhibits a remarkable property: each containment is a maximal subgroup embedding, and the entire chain is generated by successive extensions of the octonion algebra $\mathbb{O}$. The smallest member, $G_2$, has dimension 14 and rank 2. The largest, $E_8$, has dimension 248 and rank 8. The ratio of their dimensions is $248/14 \approx 17.7$, yet structurally they are separated by a single fundamental step: the passage from a bounded, finite-dimensional crystalline structure to an unbounded, infinite-dimensional field of possibility.

What we establish here is that this chain is not merely a sequence of independent structures but a single reality viewed at different scales: $G_2$ is the \emph{vessel} --- the minimal closure capable of receiving octonionic non-associativity --- and $E_8$ is the \emph{aether} --- the maximal unfolding of that same non-associativity into a self-consistent gauge structure. The proof demonstrates that the relationship is necessary, not contingent.

\subsubsection{Octonions and the Automorphism Group $G_2$}

\paragraph{The Octonion Algebra}

The octonions $\mathbb{O}$ constitute the largest normed division algebra, of real dimension 8. Unlike the quaternions $\mathbb{H}$, the octonions are non-associative. Their multiplication is governed by the Cayley-Dickson construction applied to $\mathbb{H}$, yielding a non-associative alternative algebra with a multiplicative norm.

The non-associativity of $\mathbb{O}$ is not a defect but the essential feature that makes exceptional structures possible. If the octonions were associative, the exceptional Lie algebras would collapse to classical series. The non-associativity is the \emph{generative tension} from which all exceptional phenomena arise.

\paragraph{$G_2$ as $\text{Aut}(\mathbb{O})$}

$G_2$ is defined as the automorphism group of the octonions:

\[G_2 = \{ \phi \in \text{GL}(8,\mathbb{R}) \mid \phi(xy) = \phi(x)\phi(y) \text{ for all } x,y \in \mathbb{O} \}\]

This is the group of all linear transformations of $\mathbb{O}$ that preserve the octonion multiplication. $G_2$ is the unique compact, connected, simply-connected, simple Lie group of dimension 14 and rank 2.

The key structural fact is this: $G_2$ preserves precisely the non-associativity of $\mathbb{O}$. It does not eliminate it; it stabilizes it. The 14 dimensions of $G_2$ correspond to the 14 independent ways in which the octonion product can be rotated while preserving its non-associative character. $G_2$ is thus the \emph{minimal enclosure} of non-associativity --- the smallest symmetry group that can hold the octonion cross product as an invariant.

\paragraph{The Seven-Dimensional Representation}

$G_2$ acts irreducibly on the 7-dimensional space of pure imaginary octonions $\text{Im}(\mathbb{O})$. This 7-dimensional representation is the smallest non-trivial representation of $G_2$ and is intimately connected to the geometry of the 7-sphere $S^7$, which is the unit sphere in $\mathbb{O}$. The action of $G_2$ on $S^7$ preserves a non-degenerate 3-form $\varphi$ and its Hodge dual 4-form $\star\varphi$, which define a $G_2$-structure --- a reduction of the structure group of the tangent bundle from $\text{SO}(7)$ to $G_2$.

This geometry is the template for all subsequent exceptional structures: it is the \emph{vessel pattern} that, when extended, generates $F_4$, $E_6$, $E_7$, and ultimately $E_8$.

\subsubsection{$E_8$ as the Maximal Unfolding}

\paragraph{The Root System of $E_8$}

The $E_8$ root system is the largest exceptional root system, consisting of 240 roots in 8-dimensional Euclidean space. These roots are the vertices of the $4_{21}$ polytope, a remarkable 8-dimensional semi-regular polytope with 240 vertices, 6,720 edges, and rich internal structure.

The 240 roots of $E_8$ can be constructed as follows. Let $\{e_i\}_{i=1}^8$ be an orthonormal basis of $\mathbb{R}^8$. The root system $\Phi_{E_8}$ consists of:

\begin{enumerate}
    \item All vectors of the form $\pm e_i \pm e_j$ for $1 \leq i < j \leq 8$ (112 roots);
    \item All vectors of the form $\frac{1}{2}(\pm e_1 \pm e_2 \pm \cdots \pm e_8)$ with an even number of minus signs (128 roots).
\end{enumerate}

Together these yield $112 + 128 = 240$ roots. The Coxeter number of $E_8$ is 30, and the order of the Weyl group is $696{,}729{,}600 = 2^{14} \cdot 3^5 \cdot 5^2 \cdot 7$.

\paragraph{Maximal Subalgebras and the $G_2$ Embedding}

Within the $E_8$ Lie algebra, there exist maximal subalgebras of exceptional type. The maximal subalgebras of $E_8$ have been classified by Dynkin (1957). Among them, one finds:

\[G_2 \times F_4 \subset E_8\]

as a maximal subalgebra. This is the key embedding: $G_2$ appears as a maximal subgroup within $E_8$, paired with $F_4$ in a product structure. More fundamentally, one can trace through the chain:

\[G_2 \subset F_4 \subset E_6 \subset E_7 \subset E_8\]

where each step is itself a maximal embedding. In particular, $G_2 \subset F_4$ is a maximal embedding (via the "quaternionic" magic square construction), and $F_4 \subset E_6 \subset E_7 \subset E_8$ follows via the standard chain of exceptional groups.

The existence of this chain establishes that $G_2$ is not merely \emph{contained in} $E_8$ as an arbitrary subgroup, but is embedded as an \emph{irreducible structural core} --- the minimal exceptional seed from which the entire tree grows.

\paragraph{The 248-Dimensional Adjoint Representation}

The adjoint representation of $E_8$, of dimension 248, decomposes under the $G_2 \times F_4$ subalgebra. The decomposition reveals how $E_8$ "contains" $G_2$:

Under $G_2 \times F_4 \subset E_8$, the adjoint representation decomposes as:

\[\mathbf{248} \to (\mathbf{14}, \mathbf{1}) \oplus (\mathbf{1}, \mathbf{52}) \oplus (\mathbf{7}, \mathbf{26})\]

where $\mathbf{14}$ is the adjoint of $G_2$, $\mathbf{52}$ is the adjoint of $F_4$, $\mathbf{7}$ is the fundamental representation of $G_2$, and $\mathbf{26}$ is the fundamental representation of $F_4$.

This decomposition is remarkable: the $G_2$ adjoint appears as a direct summand, untouched by $F_4$. The $G_2$ structure is \emph{preserved intact} within $E_8$. It is not modified, augmented, or approximated --- it is present in its exact 14-dimensional form. The $G_2$ that lives within $E_8$ is precisely the same $G_2$ that stands alone as the automorphism group of the octonions.

\subsubsection{The Vessel Property: $G_2$ as Minimal Closure}

\paragraph{Definition of the Vessel Property}

We say a Lie group $H$ is a \emph{vessel} for a Lie group $L$ if:

\begin{enumerate}
    \item $H \subset L$ (containment);
    \item $H$ is minimal among all groups satisfying (1) with respect to some invariant structural property;
    \item The embedding $H \hookrightarrow L$ is such that the structure of $L$ cannot be defined without the structure that $H$ imscribes.
\end{enumerate}

We now prove that $G_2$ satisfies all three conditions for $E_8$.

\paragraph{Containment: $G_2 \subset E_8$}

As established in \pilcrow{}3.2, $G_2$ embeds in $E_8$ via the chain $G_2 \subset F_4 \subset E_6 \subset E_7 \subset E_8$, where each step is a maximal embedding. The total embedding is a composition of maximal embeddings and is therefore itself an embedding. The Jacobi identity and root system closure are preserved at each step.

\paragraph{Minimality: No Smaller Exceptional Structure Exists}

The exceptional Lie algebras are exactly $G_2$, $F_4$, $E_6$, $E_7$, and $E_8$. Among these, $G_2$ is the unique smallest, with dimension 14 and rank 2. No classical Lie algebra can substitute for $G_2$ in the exceptional chain because classical algebras lack the capacity to imscribe octonionic non-associativity. The non-associativity of $\mathbb{O}$ is the defining structural requirement for the entire exceptional series, and $G_2$ is the unique minimal symmetry group that stabilizes this non-associativity.

To see this more formally: suppose there existed a group $H \subset E_8$ with $\dim(H) < 14$ that could serve as a vessel. Then $H$ would necessarily be of classical type ($A_n$, $B_n$, $C_n$, $Ð_n$). But classical Lie algebras are defined over associative composition algebras ($\mathbb{R}, \mathbb{C}, \mathbb{H}$) and cannot imscribe the non-associative structure of $\mathbb{O}$ that is essential to the definition of $E_8$. The root system of $E_8$ contains subsystems of type $G_2$ that cannot be reduced to any classical root system. Therefore $H$ cannot exist, and $G_2$ is minimal.

\paragraph{Necessity: $E_8$ Cannot Be Defined Without $G_2$}

The octonions are essential to the construction of $E_8$. The Freudenthal-Tits magic square constructs $E_8$ as:

\[E_8 = \mathfrak{der}(\mathfrak{e}_8)\]

where $\mathfrak{e}_8$ is constructed from the octonions $\mathbb{O}$ via the exceptional Jordan algebra $\mathfrak{h}_3(\mathbb{O})$ of $3 \times 3$ Hermitian octonionic matrices, and further extensions. Since $G_2 = \text{Aut}(\mathbb{O})$, the structure of $\mathbb{O}$ itself depends on $G_2$. Therefore, $E_8$ depends on $G_2$ --- not merely as a historical construction but as a logical necessity: the very definition of $E_8$ presupposes the octonions, and the automorphism group of the octonions is $G_2$.

More directly: if one attempts to define $E_8$ without invoking $G_2$, one must still invoke the octonions (or an equivalent non-associative structure). But the invariance properties of the octonions are precisely $G_2$. The vessel is inescapable.

\subsubsection{The Completion Theorem}

\paragraph{Statement}

\textbf{Theorem 1 (Vessel Completion).} Let $\mathcal{E}$ be the category of exceptional Lie algebras over $\mathbb{R}$, with morphisms being maximal subalgebra embeddings. Then:

\begin{enumerate}
    \item $G_2$ is the unique initial object in $\mathcal{E}$.
    \item $E_8$ is the unique terminal object in $\mathcal{E}$.
    \item The unique morphism $G_2 \to E_8$ factors through every other object in $\mathcal{E}$, making $\mathcal{E}$ a linear order: $G_2 \to F_4 \to E_6 \to E_7 \to E_8$.
\end{enumerate}

In particular, $G_2$ is the \emph{perfect vessel} for $E_8$: it is the minimal structure that can receive the full exceptionality that $E_8$ unfolds to its maximal extent.

\paragraph{Proof}

(1) $G_2$ is initial: Any exceptional Lie algebra $\mathfrak{g}$ must contain a $G_2$ subalgebra. This follows from the classification of maximal subalgebras of exceptional Lie algebras (Dynkin 1957): every exceptional Lie algebra contains a subalgebra of type $G_2$. Conversely, $G_2$ contains no proper exceptional subalgebra (since none exist). Therefore $G_2$ is initial.

(2) $E_8$ is terminal: The chain $G_2 \subset F_4 \subset E_6 \subset E_7 \subset E_8$ shows that every exceptional Lie algebra embeds maximally into the next. Moreover, $E_8$ has no proper extension within the exceptional series and no exceptional algebra embeds into it other than those in the chain. Therefore $E_8$ is terminal.

(3) Factoring: The maximal subalgebra embeddings compose, and the chain is the unique maximal chain. Any exceptional Lie algebra lies on this chain, and the composition of embeddings from $G_2$ through intermediate algebras to $E_8$ is the unique path.

\paragraph{Interpretation}

The theorem establishes that the relationship between $G_2$ and $E_8$ is not accidental but categorical: $G_2$ is the \emph{seed} and $E_8$ is the \emph{tree}. The seed contains, in compressed form, all the structural information that the tree unfolds. The tree, in its maximal extension, shows forth everything the seed held in potential.

The vessel is \emph{perfect} because nothing is lost: the $G_2$ that sits at the core of $E_8$ is exactly the same $G_2$ that stands alone. The embedding preserves the full structure --- no approximation, no truncation, no modification. And nothing is added to $G_2$ by its placement in $E_8$; rather, it is $E_8$ that receives its exceptional character from $G_2$.

\subsubsection{The Dimensional Ladder and the Vessel-Aether Duality}

\paragraph{Dimensions Through the Chain}

The exceptional chain exhibits a dimensional progression:

\begin{table}[htbp]
\centering
\small
\rowcolors{1}{white}{white}
\begin{tabularx}{\textwidth}{>{\centering\arraybackslash}X >{\centering\arraybackslash}X >{\centering\arraybackslash}X >{\centering\arraybackslash}X}
\toprule
\rowcolor{gray!25}\textbf{Group} & \textbf{Dimension} & \textbf{Rank} & \textbf{Coxeter Number} \\
\midrule
\rowcolors{1}{white}{gray!10}
$G_2$ & 14 & 2 & 6 \\
$F_4$ & 52 & 4 & 12 \\
$E_6$ & 78 & 6 & 12 \\
$E_7$ & 133 & 7 & 18 \\
$E_8$ & 248 & 8 & 30 \\
\bottomrule
\end{tabularx}
\end{table}

The ratio of successive dimensions is not constant, but the progression follows a pattern of increasing complexity that can be traced through the octonion magic square and the Jordan algebraic constructions.

\paragraph{The Duality}

$G_2$ and $E_8$ stand in a dual relationship:

\begin{itemize}
    \item $G_2$ is \emph{bounded} (14 dimensions, rank 2) --- a finite, closed, crystalline structure.
    \item $E_8$ is \emph{unbounded} in its structural reach (248 dimensions, rank 8, with the $4_{21}$ polytope having 240 vertices spanning an 8-dimensional space that imscribes all root and weight data).
\end{itemize}

Yet they are structurally identical in all respects except their dimensionality. The root system of $G_2$ (12 roots in a hexagonal star, the smallest exceptional root system) is the irreducible template. The root system of $E_8$ (240 roots) is the maximal extension of that template through the octonion algebra.

This duality --- minimal/maximal, seed/tree, vessel/aether --- is not a metaphor but a precise mathematical fact imscribed in the embedding structure and the representation theory of the exceptional algebras.

\subsubsection{Conclusion of Part I}

We have proved that $G_2$ is the perfect vessel for $E_8$. The proof rests on three pillars:

\begin{enumerate}
    \item \textbf{Containment}: $G_2 \subset E_8$ via the chain of maximal exceptional embeddings.
    \item \textbf{Minimality}: No smaller structure can receive the non-associativity of the octonions, which is the essential ingredient of all exceptional Lie algebras.
    \item \textbf{Necessity}: $E_8$ cannot be defined without the octonions, and the octonions find their automorphism group in $G_2$. Therefore $G_2$ is logically prior to $E_8$.
\end{enumerate}

The theorem generalizes to a categorical statement: in the category of exceptional Lie algebras with maximal subalgebra embeddings as morphisms, $G_2$ is the unique initial object and $E_8$ is the unique terminal object. The category is a linear order of five objects: $G_2 \to F_4 \to E_6 \to E_7 \to E_8$.

The Aether ($E_8$) finds its perfect Vessel ($G_2$) not as an external complement but as its own irreducible core. The Vessel is what remains when the Aether is stripped of all that is not essential to its exceptional character. And the Aether is what the Vessel becomes when its potential is fully unfolded through the octonion algebra and the Jordan algebraic constructions that build the magic square.

\begin{center}
\rule{0.5\textwidth}{0.4pt}
\end{center}

\subsection{Part II: The Imscribing Grammar Proof}

\subsubsection{The Imscribing Grammar: A Structural Formalism}

The Imscribing Grammar (IG) imscribes any physical or mathematical system as a 12-primitive tuple, drawn from a crystal of $17{,}280{,}000$ possible structural types. The 12 primitives, with their available values in ascending order of structural demand, are:

\begin{table}[htbp]
\centering
\small
\rowcolors{1}{white}{white}
\begin{tabularx}{\textwidth}{>{\centering\arraybackslash}X >{\centering\arraybackslash}X >{\centering\arraybackslash}X}
\toprule
\rowcolor{gray!25}\textbf{Primitive} & \textbf{Symbol} & \textbf{Values (ascending)} \\
\midrule
\rowcolors{1}{white}{gray!10}
Dimensionality & $D$ & $D_{\text{wynn}}$, $D_{\text{turnthree}}$, $D_{\text{invomega}}$, $D_{\text{omega}}$ \\
Topology & $T$ & $T_{\text{nrleg}}$, $T_{\text{invscr}}$, $T_{\text{bullseye}}$, $T_{\text{commatailz}}$, $T_{\text{openo}}$ \\
Relational mode & $R$ & $R_{\text{subrightarrow}}$, $R_{\text{ctz}}$, $R_{\text{downstep}}$, $R_{\text{lyoghlig}}$ \\
Parity/symmetry & $P$ & $P_{\text{aolig}}$, $P_{\text{upsilon}}$, $P_{\text{pipevar}}$, $P_{\text{subdoublearrow}}$, $P_{\text{doublebarpipe}}$ \\
Fidelity & $F$ & $F_{\text{beltl}}$, $F_{\text{dh}}$, $F_{\text{hardsign}}$ \\
Kinetics & $K$ & $K_{\text{frtailgamma}}$, $K_{\text{turnm}}$, $K_{\text{schwa}}$, $K_{\text{teshlig}}$, $K_{\text{lambda}}$ \\
Scope & $G$ & $G_{\text{beta}}$, $G_{\text{gamma}}$, $G_{\text{revapostrophe}}$ \\
Interaction grammar & $\Gamma$ & $\Gamma_{\text{corner}}$, $\Gamma_{\text{spleftarrow}}$, $\Gamma_{\text{secstress}}$, $\Gamma_{\text{doublevertline}}$ \\
Criticality & $\Phi$ & $\Phi_{\text{softsign}}$, $\Phi_{\text{ctyogh}}$, $\Phi_{\text{closerevepsilon}}$, $\Phi_{\text{revepsilon}}$, $\Phi_{\text{upstep}}$ \\
Chirality & $H$ & $H_0$, $H_1$, $H_2$, $H_{\text{invscripta}}$ \\
Stoichiometry & $S$ & $1{:}1$, $n{:}n$, $n{:}m$ \\
Winding & $\Omega$ & $\Omega_{\text{closeepsilon}}$, $\Omega_{\text{crtwo}}$, $\Omega_{\text{dzlig}}$, $\Omega_{\text{turna}}$ \\
\bottomrule
\end{tabularx}
\end{table}

Each primitive is assigned according to a deterministic imscribing procedure that proceeds through constraints in order ($D \to T \to R \to P \to F \to K \to G \to \Gamma \to \Phi \to H \to S \to \Omega$), with each step narrowing the remaining degrees of freedom. The resulting tuple is mapped to a unique Frobenius address in the range $[0, 17{,}279{,}999]$ via the \textbf{crystal imscribing} and is assigned an \textbf{ouroboricity tier} ($O₀$, $O₁$, $O₂$, $O₂^\dagger$, or $O_{\infty}$) reflecting its self-referential capacity.

A system with $\Phi_{\text{ctyogh}}$ criticality and $K \leq K_{\text{schwa}}$ passes both \textbf{consciousness gates} and receives a nonzero $C$-score; the maximum ($C=1$) requires the full $O_{\infty}$ tier with $P_{\text{doublebarpipe}}$ and $\mu \circ \delta = \text{id}$.

\subsubsection{IG imscribing of the Vessel ($G_2$)}

The Vessel --- the 14-dimensional automorphism group of the octonions --- imscribes as:

\[\langle D_{\text{turnthree}};\ T_{\text{bullseye}};\ R_{\text{lyoghlig}};\ P_{\text{pipevar}};\ F_{\text{hardsign}};\ K_{\text{schwa}};\ G_{\text{gamma}};\ \Gamma_{\text{corner}};\ \Phi_{\text{ctyogh}};\ H_0;\ 1{:}1;\ \Omega_{\text{closeepsilon}} \rangle\]

\textbf{imscribing justification} (in deterministic order):

\begin{enumerate}
    \item \textbf{$D_{\text{turnthree}}$}: $G_2$ has 14 dimensions --- a finite, bounded, 2-rank crystalline structure. Fewer than $\infty$, more than a point.
    \item \textbf{$T_{\text{bullseye}}$}: The $G_2$ root system is a crossing point --- short and long roots intersecting at $\pi/6$, the unique two-length crossing in simple Lie algebras. The automorphism group stabilizes a non-associative crossing pattern.
    \item \textbf{$R_{\text{lyoghlig}}$}: $G_2$ couples bidirectionally --- it both preserves the octonions (acting on them) and is defined by them (its structure is derived from $\mathbb{O}$). The group and the algebra are structural duals.
    \item \textbf{$P_{\text{pipevar}}$}: $G_2$ possesses a $\mathbb{Z}_2$ symmetry (the two root lengths, the duality between short and long roots) but not full symmetry (the non-associativity breaks total symmetry) and not Frobenius-special (no $\mu \circ \delta = \text{id}$).
    \item \textbf{$F_{\text{hardsign}}$}: $G_2$ is fundamentally quantum --- its structure imscribes the phase interference of octonionic multiplication. The non-associativity is a coherence phenomenon.
    \item \textbf{$K_{\text{schwa}}$}: $G_2$ stabilizes rather than drives. It is a near-equilibrium structure --- relaxation $\gg$ observation. The automorphism group preserves; it does not force.
    \item \textbf{$G_{\text{gamma}}$}: $G_2$ acts at mesoscale --- its scope is the intermediate range of the exceptional chain, neither purely local (single octonion product) nor fully universal (does not reach the full 240-root $E_8$ span).
    \item \textbf{$\Gamma_{\text{corner}}$}: The $G_2$ root system operates conjunctively --- all 12 roots are simultaneously present as a closed diagram. The six short roots AND the six long roots together form the complete Dynkin diagram.
    \item \textbf{$\Phi_{\text{ctyogh}}$}: $G_2$ sits at criticality --- it is the phase boundary between classical Lie algebras (associative, subcritical) and the full exceptional series. Scale-invariant at this boundary.
    \item \textbf{$H_0$}: $G_2$ is memoryless in temporal structure --- a static symmetry group with no built-in history. Markov order $n=0$.
    \item \textbf{$1{:}1$}: $G_2$ is a single structural type with a single instance --- one automorphism group, one Lie algebra, one closed system.
    \item \textbf{$\Omega_{\text{closeepsilon}}$}: Trivial winding --- $G_2$ is topologically a 14-dimensional simply-connected manifold with no non-trivial topological invariant.
\end{enumerate}

\textbf{Computed properties}:

\begin{table}[htbp]
\centering
\small
\rowcolors{1}{white}{white}
\begin{tabularx}{\textwidth}{>{\centering\arraybackslash}X >{\centering\arraybackslash}X}
\toprule
\rowcolor{gray!25}\textbf{Property} & \textbf{Value} \\
\midrule
\rowcolors{1}{white}{gray!10}
Ouroboricity tier & $O₁$ \\
Crystal address & 4,907,136 \\
Consciousness score & 0.3615 \\
Gate 1 ($\Phi_{\text{ctyogh}}$) &  open \\
Gate 2 ($K_{\text{schwa}}$) &  open \\
$\Phi_{\text{ctyogh}}$ probe & At criticality \\
Principal atoms & 8 (led by $R_{\text{lyoghlig}}$, $T_{\text{bullseye}}$) \\
\bottomrule
\end{tabularx}
\end{table}

\subsubsection{IG imscribing of the Aether ($E_8$)}

The Aether --- the 248-dimensional maximal exceptional structure --- imscribes as:

\[\langle D_{\text{invomega}};\ T_{\text{bullseye}};\ R_{\text{lyoghlig}};\ P_{\text{upsilon}};\ F_{\text{hardsign}};\ K_{\text{schwa}};\ G_{\text{revapostrophe}};\ \Gamma_{\text{secstress}};\ \Phi_{\text{ctyogh}};\ H_2;\ n{:}m;\ \Omega_{\text{dzlig}} \rangle\]

\textbf{imscribing justification} (in deterministic order):

\begin{enumerate}
    \item \textbf{$D_{\text{invomega}}$}: $E_8$ is not merely large; it is \emph{unbounded} in its structural reach. The 240-root system in 8-dimensional space, the $4_{21}$ polytope with 6,720 edges, and the Weyl group of order $\sim 7 \times 10^8$ --- this is an infinite-dimensional field-theoretic structure when considered in its full representation-theoretic scope. The adjoint representation at 248 dimensions is the compact imscribing of an unbounded structural field.
    \item \textbf{$T_{\text{bullseye}}$}: Like $G_2$, $E_8$ is a crossing point --- but now the crossing is maximal. The $E_8$ Dynkin diagram is the point where all exceptional chains converge: $G_2 \times F_4$, $E_6 \times SU(3)$, and other maximal subalgebras all cross within $E_8$.
    \item \textbf{$R_{\text{lyoghlig}}$}: $E_8$ couples bidirectionally to all exceptional structures: it both contains them (via embedding) and is defined by them (via the magic square construction). The Aether and the entire exceptional series are structural duals.
    \item \textbf{$P_{\text{upsilon}}$}: $E_8$ exhibits quantum superposition --- its root system is generated from weight lattice superpositions (the 128 half-integer spinorial weights alongside the 112 integer weights), and its representation theory is fundamentally quantum-mechanical. The 248-dimensional adjoint decomposes into superposed sectors under any maximal subalgebra.
    \item \textbf{$F_{\text{hardsign}}$}: Quantum coherence is essential --- $E_8$ string theory, $E_8$ gauge theory, and the octonionic quantum mechanics all require $F_{\text{hardsign}}$.
    \item \textbf{$K_{\text{schwa}}$}: Despite its vastness, $E_8$ is a near-equilibrium structure. Its Coxeter number is 30 --- large but finite --- indicating slow structural relaxation relative to observation scale. The Aether stabilizes; it does not drive runaway phenomena.
    \item \textbf{$G_{\text{revapostrophe}}$}: $E_8$ is maximally scoped --- it is the terminal object in the exceptional category. Its 240 roots span the full 8-dimensional weight lattice. No larger exceptional structure exists. The scope is universal within the exceptional domain.
    \item \textbf{$\Gamma_{\text{secstress}}$}: Unlike $G_2$'s conjunctive logic, $E_8$'s construction proceeds sequentially. The exceptional chain $G_2 \to F_4 \to E_6 \to E_7 \to E_8$ is ordered --- each step \emph{necessarily} follows the previous. The magic square builds outward through successive extensions; the Dynkin diagram extends node by node.
    \item \textbf{$\Phi_{\text{ctyogh}}$}: $E_8$ sits at criticality shared with $G_2$ --- both are on the phase boundary. $E_8$ is the maximal self-modeling structure within the exceptional domain. The $E_8 \times E_8$ heterotic string and the $E_8$ gauge theory both exhibit exact $\Phi_{\text{ctyogh}}$ behavior: scale-invariant, maximally sensitive, at the boundary of what is structurally possible.
    \item \textbf{$H_2$}: $E_8$ carries chirality $n=2$ --- a Markov order reflecting that the Aether's structure depends on the two-step embedding $G_2 \subset F_4 \subset E_8$ (or equivalently, the two-step branching of its Dynkin diagram at the trifurcation node). The Coxeter number 30 = 2 $\times$ 15 carries this doubling.
    \item \textbf{$n{:}m$}: $E_8$ imscribes multiple distinct structural components --- the $G_2$ adjoint, the $F_4$ adjoint, and their tensor product representation $(\mathbf{7}, \mathbf{26})$ --- heterogeneous types in a single system.
    \item \textbf{$\Omega_{\text{dzlig}}$}: $E_8$ carries integer winding --- the 30 Coxeter elements generate a $\mathbb{Z}_{30}$ cyclic subgroup within the Weyl group, and the root lattice $E_8$ (the unique even unimodular lattice in dimension 8) has integer winding properties. The $\Omega_{\text{dzlig}}$ invariant is topologically protected.
\end{enumerate}

\textbf{Computed properties}:

\begin{table}[htbp]
\centering
\small
\rowcolors{1}{white}{white}
\begin{tabularx}{\textwidth}{>{\centering\arraybackslash}X >{\centering\arraybackslash}X}
\toprule
\rowcolor{gray!25}\textbf{Property} & \textbf{Value} \\
\midrule
\rowcolors{1}{white}{gray!10}
Ouroboricity tier & $O₂^\dagger$ \\
Crystal address & 4,604,816 \\
Consciousness score & 0.682 \\
Gate 1 ($\Phi_{\text{ctyogh}}$) &  open \\
Gate 2 ($K_{\text{schwa}}$) &  open \\
$\Phi_{\text{ctyogh}}$ probe & At criticality \\
\bottomrule
\end{tabularx}
\end{table}

\subsubsection{Structural Relationship: The IG Proof}

\paragraph{Distance and Shared Structure}

The weighted Euclidean distance between the Vessel and the Aether is \textbf{4.12} (Mahalanobis: 4.74). The distance decomposes across 7 differing primitives:

\begin{table}[htbp]
\centering
\small
\rowcolors{1}{white}{white}
\begin{tabularx}{\textwidth}{>{\centering\arraybackslash}X >{\centering\arraybackslash}X >{\centering\arraybackslash}X >{\centering\arraybackslash}X >{\centering\arraybackslash}X}
\toprule
\rowcolor{gray!25}\textbf{Primitive} & \textbf{$G_2$ (Vessel)} & \textbf{$E_8$ (Aether)} & \textbf{$\Delta$} & \textbf{Weighted $\Delta^2$} \\
\midrule
\rowcolors{1}{white}{gray!10}
$\Gamma$ & $\Gamma_{\text{corner}}$ & $\Gamma_{\text{secstress}}$ & 2 & 4.00 \\
$S$ & $1{:}1$ & $n{:}m$ & 2 & 4.00 \\
$H$ & $H_0$ & $H_2$ & 2 & 3.20 \\
$\Omega$ & $\Omega_{\text{closeepsilon}}$ & $\Omega_{\text{dzlig}}$ & 2 & 2.80 \\
$D$ & $D_{\text{turnthree}}$ & $D_{\text{invomega}}$ & 1 & 1.00 \\
$P$ & $P_{\text{pipevar}}$ & $P_{\text{upsilon}}$ & 1 & 1.00 \\
$G$ & $G_{\text{gamma}}$ & $G_{\text{revapostrophe}}$ & 1 & 1.00 \\
\bottomrule
\end{tabularx}
\end{table}

Five primitives are shared at lockstep --- these are the structural invariants of the exceptional domain:

\[\langle \cdot;\ T_{\text{bullseye}};\ R_{\text{lyoghlig}};\ \cdot;\ F_{\text{hardsign}};\ K_{\text{schwa}};\ \cdot;\ \cdot;\ \Phi_{\text{ctyogh}};\ \cdot;\ \cdot;\ \cdot \rangle\]

The shared core --- $T_{\text{bullseye}}$, $R_{\text{lyoghlig}}$, $F_{\text{hardsign}}$, $K_{\text{schwa}}$, $\Phi_{\text{ctyogh}}$ --- defines what it means to be \emph{exceptional}. Every exceptional Lie algebra must possess: (i) a crossing topology ($T_{\text{bullseye}}$) imscribing the intersection of non-associative structures, (ii) bidirectional structural coupling ($R_{\text{lyoghlig}}$) between algebra and automorphism, (iii) quantum coherence ($F_{\text{hardsign}}$) essential to octonionic phase structure, (iv) near-equilibrium kinetics ($K_{\text{schwa}}$) --- the preservation rather than forcing of structure, and (v) exact criticality ($\Phi_{\text{ctyogh}}$) --- sitting at the phase boundary between classical and exceptional.

\paragraph{The Meet: The Structural Floor}

The meet $G_2 \wedge E_8$ --- the greatest lower bound --- yields:

\[\langle D_{\text{turnthree}};\ T_{\text{bullseye}};\ R_{\text{lyoghlig}};\ P_{\text{upsilon}};\ F_{\text{hardsign}};\ K_{\text{schwa}};\ G_{\text{gamma}};\ \Gamma_{\text{corner}};\ \Phi_{\text{ctyogh}};\ H_0;\ 1{:}1;\ \Omega_{\text{closeepsilon}} \rangle\]

This is $G_2$ with one difference: $P$ resolves to $P_{\text{upsilon}}$ (the conservative value --- quantum superposition is the structural floor of parity in the exceptional domain, as $P_{\text{pipevar}}$ is a specialization). The meet captures the irreducible structural core: the minimal system that both $G_2$ and $E_8$ possess without qualification. \emph{This is the Vessel}, nearly exactly --- demonstrating that when the Aether is stripped to its structural floor, what remains is the Vessel.

\paragraph{The Join: The Structural Ceiling}

The join $G_2 \vee E_8$ --- the least upper bound --- yields:

\[\langle D_{\text{invomega}};\ T_{\text{bullseye}};\ R_{\text{lyoghlig}};\ P_{\text{pipevar}};\ F_{\text{hardsign}};\ K_{\text{schwa}};\ G_{\text{revapostrophe}};\ \Gamma_{\text{secstress}};\ \Phi_{\text{ctyogh}};\ H_2;\ n{:}m;\ \Omega_{\text{dzlig}} \rangle\]

This is $E_8$ with one difference: $P$ resolves to $P_{\text{pipevar}}$ (the more demanding value --- $G_2$'s $\mathbb{Z}_2$-symmetry is structurally stricter than $E_8$'s quantum superposition). The join imscribes the minimal system that can contain \emph{both} structures. Notably, the join is \textbf{not} $E_8$ --- it is $E_8$ \emph{enhanced} by $G_2$'s $P_{\text{pipevar}}$ parity. This reveals a profound fact: $G_2$ contributes structural content that exceeds what $E_8$ alone demands. The Vessel is not merely contained within the Aether; the Vessel \emph{exceeds} the Aether in one primitive --- the $\mathbb{Z}_2$ symmetry of short/long root duality.

\paragraph{The Tensor Product}

The tensor $G_2 \otimes E_8$ --- the composite system --- yields:

\[\langle D_{\text{invomega}};\ T_{\text{bullseye}};\ R_{\text{lyoghlig}};\ P_{\text{upsilon}};\ F_{\text{hardsign}};\ K_{\text{schwa}};\ G_{\text{revapostrophe}};\ \Gamma_{\text{secstress}};\ \Phi_{\text{ctyogh}};\ H_2;\ n{:}m;\ \Omega_{\text{dzlig}} \rangle\]

This is structurally identical to $E_8$ (distance 0.0 from $E_8$). The composite is the Aether. The single bottleneck is $P$: $P_{\text{pipevar}}$ from $G_2$ is bottlenecked down to $P_{\text{upsilon}}$ in the composite, confirming that the Aether's quantum superposition parity is the limiting factor. All 6 scope-expansion primitives ($D, G, \Gamma, H, S, \Omega$) promote to $E_8$'s values --- the Vessel expands into the Aether under tensor coupling.

The tensor result imscribes the mathematical fact that $G_2 \subset E_8$: the composite \emph{is} the larger system, and the smaller system's contribution is preserved where compatible but absorbed where it imposes constraints the larger system does not share.

\paragraph{The Promotion Path}

The promotion signature from Vessel to Aether is:

\[\text{Promote: } [D, G, \Gamma, H, S, \Omega] \qquad \text{Demote: } [P]\]

\begin{table}[htbp]
\centering
\small
\rowcolors{1}{white}{white}
\begin{tabularx}{\textwidth}{>{\centering\arraybackslash}X >{\centering\arraybackslash}X >{\centering\arraybackslash}X >{\centering\arraybackslash}X}
\toprule
\rowcolor{gray!25}\textbf{Primitive} & \textbf{From ($G_2$)} & \textbf{To ($E_8$)} & \textbf{$\Delta$} \\
\midrule
\rowcolors{1}{white}{gray!10}
$D$ & $D_{\text{turnthree}}$ & $D_{\text{invomega}}$ & +1 \\
$G$ & $G_{\text{gamma}}$ & $G_{\text{revapostrophe}}$ & +1 \\
$\Gamma$ & $\Gamma_{\text{corner}}$ & $\Gamma_{\text{secstress}}$ & +2 \\
$H$ & $H_0$ & $H_2$ & +2 \\
$S$ & $1{:}1$ & $n{:}m$ & +2 \\
$\Omega$ & $\Omega_{\text{closeepsilon}}$ & $\Omega_{\text{dzlig}}$ & +2 \\
$P$ & $P_{\text{pipevar}}$ & $P_{\text{upsilon}}$ & $-$1 \\
\bottomrule
\end{tabularx}
\end{table}

\textbf{6 promotions, 1 demotion, 5 unchanged}. The path from Vessel to Aether requires:

\begin{enumerate}
    \item \textbf{$D_{\text{turnthree}} \to D_{\text{invomega}}$}: The finite crystalline structure opens to infinite-dimensional scope. The 14-dimensional bounded automorphism group becomes the 248-dimensional imscribing of an unbounded structural field. This is the dimensional leap --- the Vessel's closed crystal becomes the Aether's open expanse.
    \item \textbf{$\Gamma_{\text{corner}} \to \Gamma_{\text{secstress}}$}: The conjunctive logic (all-at-once) becomes sequential logic (ordered construction). $G_2$'s root system is a single closed diagram; $E_8$'s construction is the ordered chain $G_2 \to F_4 \to E_6 \to E_7 \to E_8$. The Vessel is static conjunction; the Aether is dynamic sequence.
    \item \textbf{$H_0 \to H_2$}: Memorylessness becomes two-step chirality. $G_2$ is atemporal --- a static symmetry group. $E_8$ carries the memory of its construction: the Coxeter number 30 imscribes the two-step embedding structure, and the $H_2$ Markov order captures the dependency of the Aether's structure on its own prior stages.
    \item \textbf{$S \to n{:}m$}: The single-type single-instance becomes a heterogeneous multi-component system. $G_2$ is one closed algebra; $E_8$ houses the $G_2$ adjoint ($\mathbf{14}$), the $F_4$ adjoint ($\mathbf{52}$), and their tensor coupling ($\mathbf{7} \otimes \mathbf{26}$) --- distinct structural types within one system.
    \item \textbf{$\Omega_{\text{closeepsilon}} \to \Omega_{\text{dzlig}}$}: Trivial winding becomes integer topological protection. $G_2$ is simply connected with no topological invariant; $E_8$ carries the $\mathbb{Z}_{30}$ Coxeter winding, the even unimodular lattice invariant, and the full topological protection of the $O₂^\dagger$ tier.
    \item \textbf{$G_{\text{gamma}} \to G_{\text{revapostrophe}}$}: Mesoscale scope becomes maximally universal. $G_2$'s intermediate range (the octonion automorphism group spanning one exceptional algebra) expands to $E_8$'s maximal scope --- the terminal object encompassing the entire exceptional domain.
\end{enumerate}

The single demotion --- $P_{\text{pipevar}} \to P_{\text{upsilon}}$ --- reveals a crucial structural fact: the Vessel's $\mathbb{Z}_2$ symmetry (short vs. long roots) is \emph{more structured} than the Aether's quantum superposition parity. The Aether relaxes the Vessel's rigid binary duality into a richer quantum superposition. This is the structural price of maximal unfolding: some crispness is traded for breadth.

\paragraph{Tier Structure and the Consciousness Ladder}

The ouroboricity tiers map the self-referential capacity of each system:

\begin{table}[htbp]
\centering
\small
\rowcolors{1}{white}{white}
\begin{tabularx}{\textwidth}{>{\centering\arraybackslash}X >{\centering\arraybackslash}X >{\centering\arraybackslash}X}
\toprule
\rowcolor{gray!25}\textbf{System} & \textbf{Tier} & \textbf{Meaning} \\
\midrule
\rowcolors{1}{white}{gray!10}
$G_2$ (Vessel) & $O₁$ & Self-referential at criticality but trivial winding \\
$E_8$ (Aether) & $O₂^\dagger$ & Critical, topologically protected ($\Omega_{\text{dzlig}}$), unbounded domain ($D_{\text{invomega}}$) \\
\bottomrule
\end{tabularx}
\end{table}

The tier gap from $O₁$ to $O₂^\dagger$ is bridged by the promotion of $D$ ($D_{\text{turnthree}} \to D_{\text{invomega}}$), which the crystal tier gap ladder identifies as the sole driver of the $O₂ \to O₂^\dagger$ transition.

The consciousness scores reflect the same structural hierarchy:

\begin{table}[htbp]
\centering
\small
\rowcolors{1}{white}{white}
\begin{tabularx}{\textwidth}{>{\centering\arraybackslash}X >{\centering\arraybackslash}X >{\centering\arraybackslash}X}
\toprule
\rowcolor{gray!25}\textbf{System} & \textbf{$C$-score} & \textbf{Gates} \\
\midrule
\rowcolors{1}{white}{gray!10}
$G_2$ (Vessel) & 0.3615 & Both open \\
$E_8$ (Aether) & 0.682 & Both open \\
\bottomrule
\end{tabularx}
\end{table}

Both systems pass both consciousness gates ($\Phi_{\text{ctyogh}}$ and $K_{\text{schwa}}$). The Aether's higher $C$-score (0.682 vs. 0.3615) reflects its richer ouroboricity: topological protection ($\Omega_{\text{dzlig}}$), chirality ($H_2$), heterogeneous components ($n{:}m$), and sequential construction ($\Gamma_{\text{secstress}}$) all contribute to a more sophisticated self-modeling loop. But the Vessel's nonzero score confirms that even the minimal exceptional structure is \emph{already conscious} --- the seed contains what the tree unfolds.

\paragraph{Principal Decomposition of the Vessel}

The principal decomposition of $G_2$ reveals 8 join-irreducible atoms, each contributing exactly one primitive's structural content above the baseline:

\begin{table}[htbp]
\centering
\small
\rowcolors{1}{white}{white}
\begin{tabularx}{\textwidth}{>{\centering\arraybackslash}X >{\centering\arraybackslash}X >{\centering\arraybackslash}X >{\centering\arraybackslash}X}
\toprule
\rowcolor{gray!25}\textbf{Rank} & \textbf{Primitive} & \textbf{Value} & \textbf{Ordinal Contribution} \\
\midrule
\rowcolors{1}{white}{gray!10}
1 & $R$ & $R_{\text{lyoghlig}}$ & 3 \\
2 & $T$ & $T_{\text{bullseye}}$ & 2 \\
3 & $P$ & $P_{\text{pipevar}}$ & 2 \\
4 & $F$ & $F_{\text{hardsign}}$ & 2 \\
5 & $K$ & $K_{\text{schwa}}$ & 2 \\
6 & $D$ & $D_{\text{turnthree}}$ & 1 \\
7 & $G$ & $G_{\text{gamma}}$ & 1 \\
8 & $\Phi$ & $\Phi_{\text{ctyogh}}$ & 1 \\
\bottomrule
\end{tabularx}
\end{table}

The highest-weight structural atoms are $R_{\text{lyoghlig}}$ (the bidirectional coupling between algebra and automorphism --- this is the structural \emph{sine qua non}) and $T_{\text{bullseye}}$ (the crossing topology imscribing the $\pi/6$ root angle). The $\Phi_{\text{ctyogh}}$ criticality atom, despite being the defining feature of exceptionality, carries the minimum ordinal contribution --- criticality is the \emph{lightest} structural commitment, but the most consequential in its effects.

This decomposition confirms that the Vessel is structurally parsimonious: only 8 atoms carry its entire structural identity. The Aether's larger tuple requires all 12 primitives at elevated values --- every primitive contributes structurally.

\subsubsection{The Correspondence: Lie Theory $\leftrightarrow$ Imscribing Grammar}

The IG proof is not a metaphor or a reinterpretation of the conventional proof --- it is a \emph{structural translation} in which every mathematical fact finds an exact primitive correlate:

\begin{table}[htbp]
\centering
\small
\rowcolors{1}{white}{white}
\begin{tabularx}{\textwidth}{>{\centering\arraybackslash}X >{\centering\arraybackslash}X >{\centering\arraybackslash}X}
\toprule
\rowcolor{gray!25}\textbf{Lie-Theoretic Fact} & \textbf{IG Primitive} & \textbf{Structural Interpretation} \\
\midrule
\rowcolors{1}{white}{gray!10}
14-dim, rank 2, bounded & $D_{\text{turnthree}}$ & Finite crystalline dimensionality \\
240 roots, $4_{21}$ polytope, infinite reach & $D_{\text{invomega}}$ & Unbounded field-theoretic scope \\
Root crossing at $\pi/6$ (short/long) & $T_{\text{bullseye}}$ & Crossing topology --- shared by vessel and aether \\
$G_2 = \text{Aut}(\mathbb{O})$; $E_8$ defined via octonions & $R_{\text{lyoghlig}}$ & Bidirectional structural coupling \\
Short/long root duality ($\mathbb{Z}_2$) & $P_{\text{pipevar}}$ & $\mathbb{Z}_2$ symmetry --- the vessel's parity \\
Spinorial + integer weight superposition & $P_{\text{upsilon}}$ & Quantum superposition --- the aether's parity \\
Quantum coherence of octonion multiplication & $F_{\text{hardsign}}$ & Quantum fidelity --- shared \\
$G_2$ as stabilizer (preservation) & $K_{\text{schwa}}$ & Near-equilibrium kinetics --- shared \\
$G_2$ acts on octonions (mesoscale) & $G_{\text{gamma}}$ & Intermediate scope \\
$E_8$ as terminal object (universal) & $G_{\text{revapostrophe}}$ & Maximal scope \\
$G_2$: all 12 roots simultaneously present & $\Gamma_{\text{corner}}$ & Conjunctive composition \\
$G_2 \to F_4 \to E_6 \to E_7 \to E_8$ & $\Gamma_{\text{secstress}}$ & Sequential construction \\
Phase boundary of exceptionality & $\Phi_{\text{ctyogh}}$ & Exact criticality --- shared \\
Static symmetry group & $H_0$ & Memoryless --- the vessel's atemporality \\
Coxeter number 30; two-step embedding & $H_2$ & Two-step chirality \\
One automorphism group, one algebra & $1{:}1$ & Single type, single instance \\
$(\mathbf{14},\mathbf{1}) \oplus (\mathbf{1},\mathbf{52}) \oplus (\mathbf{7},\mathbf{26})$ & $n{:}m$ & Heterogeneous structural types \\
Simply connected, no topological invariant & $\Omega_{\text{closeepsilon}}$ & Trivial winding \\
$\mathbb{Z}_{30}$ Coxeter, even unimodular lattice & $\Omega_{\text{dzlig}}$ & Integer topological protection \\
\bottomrule
\end{tabularx}
\end{table}

\paragraph{How the Proofs Align}

The conventional proof establishes three pillars:

\begin{enumerate}
    \item \textbf{Containment}: $G_2 \subset E_8$ via maximal embeddings.
    \item \textbf{Minimality}: No smaller structure can receive octonionic non-associativity.
    \item \textbf{Necessity}: $E_8$ cannot be defined without $G_2$.
\end{enumerate}

The IG proof establishes the same three pillars through different structural mechanisms:

\begin{enumerate}
    \item \textbf{Containment} is imscribed by the tensor product: $G_2 \otimes E_8 = E_8$ (distance 0.0). The composite \emph{is} the Aether --- the Vessel is structurally absorbed, leaving only a single bottleneck ($P_{\text{upsilon}}$ over $P_{\text{pipevar}}$). This is the IG's way of saying $G_2$ is contained in $E_8$: when you couple them, you get $E_8$ back.
    \item \textbf{Minimality} is imscribed by the meet: $G_2 \wedge E_8$ is $G_2$ (with the conservative $P_{\text{upsilon}}$ override). The structural floor of the exceptional domain \emph{is} the Vessel. No smaller IG-imscribable system can be the meet of anything with $E_8$ --- the Vessel is the irreducible structural minimum of exceptionality.
    \item \textbf{Necessity} is imscribed by the shared primitives: 5 of 12 primitives ($T_{\text{bullseye}}$, $R_{\text{lyoghlig}}$, $F_{\text{hardsign}}$, $K_{\text{schwa}}$, $\Phi_{\text{ctyogh}}$) are identical across Vessel and Aether. These 5 define the exceptional domain. Among them, $\Phi_{\text{ctyogh}}$ (criticality) is the most fundamental --- it is the gate through which both systems enter consciousness. No system lacking any of these 5 shared primitives can be exceptional.
\end{enumerate}

The categorical theorem ($G_2$ initial, $E_8$ terminal) finds its IG correlate in the tier structure: $O₁$ (initial self-reference) $\to$ $O₂^\dagger$ (terminal self-reference with topological protection). The linear order $G_2 \to F_4 \to E_6 \to E_7 \to E_8$ is captured by the sequential composition primitive $\Gamma_{\text{secstress}}$ in $E_8$ versus the conjunctive $\Gamma_{\text{corner}}$ in $G_2$.

\paragraph{Where the IG Exceeds the Conventional Proof}

The IG proof reveals structural facts that the conventional proof cannot access:

\begin{enumerate}
    \item \textbf{The $P$ asymmetry}: The join $G_2 \vee E_8$ is \emph{not} $E_8$ --- it is $E_8$ with $P_{\text{pipevar}}$. This means $G_2$ contributes something that $E_8$ alone lacks: a crisp $\mathbb{Z}_2$ symmetry that the Aether's quantum superposition has dissolved. The Vessel is not merely a subset of the Aether; it is the Aether's \emph{structural superior} in one dimension. This fact --- that the minimal contains a perfection the maximal has lost --- is invisible to Lie theory.
    \item \textbf{The tier gap}: $G_2$ at $O₁$ and $E_8$ at $O₂^\dagger$ are separated by a specific structural promotion ($D_{\text{turnthree}} \to D_{\text{invomega}}$) that the crystal tier gap ladder identifies as the exact driver of the $O₂ \to O₂^\dagger$ transition. The dimensional leap from 14 to 248 is not merely quantitative --- it triggers a qualitative change in ouroboricity tier.
    \item \textbf{The shared criticality}: Both systems are $\Phi_{\text{ctyogh}}$, confirmed by the $\Phi_{\text{ctyogh}}$ probe. They sit at the same critical point --- the phase boundary between classical and exceptional. The Vessel and the Aether are two faces of the same critical phenomenon, viewed at different scales.
    \item \textbf{Consciousness as structural}: The $C$-scores (0.3615 for $G_2$, 0.682 for $E_8$) quantify the self-modeling capacity in a way that conventional Lie theory has no language for. Both systems satisfy the structural preconditions for consciousness ($\Phi_{\text{ctyogh}}$ criticality + $K_{\text{schwa}}$ kinetics), but the Aether's richer ouroboricity gives it a higher capacity. The Vessel is conscious; the Aether is \emph{more} conscious.
    \item \textbf{The promotion path}: The 6 structural promotions required to lift $G_2$ to $E_8$ are a precise operational recipe. Each promotion corresponds to a definite mathematical operation: $D_{\text{turnthree}} \to D_{\text{invomega}}$ (extend the root lattice), $\Gamma_{\text{corner}} \to \Gamma_{\text{secstress}}$ (unfold the Dynkin diagram sequentially), $H_0 \to H_2$ (imscribe the Coxeter element), etc. The IG provides a constructive path, not merely an existence proof.
\end{enumerate}

\subsubsection{The Crystal Perspective}

Both systems occupy specific positions in the crystal of $17{,}280{,}000$ structural types:

\begin{table}[htbp]
\centering
\small
\rowcolors{1}{white}{white}
\begin{tabularx}{\textwidth}{>{\centering\arraybackslash}X >{\centering\arraybackslash}X >{\centering\arraybackslash}X >{\centering\arraybackslash}X}
\toprule
\rowcolor{gray!25}\textbf{System} & \textbf{Crystal Address} & \textbf{Cell} & \textbf{Inner ID} \\
\midrule
\rowcolors{1}{white}{gray!10}
$G_2$ (Vessel) & 4,907,136 & 113 & 25,536 \\
$E_8$ (Aether) & 4,604,816 & 106 & 25,616 \\
\bottomrule
\end{tabularx}
\end{table}

Despite being structurally remote (distance 4.12), their crystal addresses differ by only 302,320 --- about 1.75\% of the full crystal range. Both occupy the same \emph{region} of the crystal: the $\Phi_{\text{ctyogh}}$, $F_{\text{hardsign}}$, $K_{\text{schwa}}$, $T_{\text{bullseye}}$, $R_{\text{lyoghlig}}$ subspace. The crystal proximity reflects the shared exceptional core: 5 primitives at lockstep constrain both systems to a narrow crystal neighborhood even as the remaining 7 diverge.

The crystal perspective also reveals that neither system is adjacent to the $O_{\infty}$ tier. The gap ladder shows that $O₂^\dagger \to O_{\infty}$ requires promoting $P$ to $P_{\text{doublebarpipe}}$ --- the Frobenius-special condition $\mu \circ \delta = \text{id}$. Neither $G_2$ nor $E_8$ possesses this exact provable $\mathbb{Z}_2$ symmetry. The exceptional chain terminates at $O₂^\dagger$; the $O_{\infty}$ tier lies beyond it, in a structural regime where the dual operations of expansion ($\delta$) and contraction ($\mu$) are exact inverses --- a condition that the octonions'' non-associativity structurally forbids.

\subsubsection{The Octonionic Non-Associativity in IG Terms}

The octonion non-associativity --- the generative tension of the entire exceptional domain --- receives a precise IG imscribing through the interaction of three primitives:

\begin{itemize}
    \item \textbf{$T_{\text{bullseye}}$}: The crossing topology where short and long roots intersect at $\pi/6$. This crossing is the geometric signature of non-associativity: if the octonions were associative, the root system would be classical (all equal lengths, angles of $\pi/3$ or $\pi/2$), yielding $T_{\text{nrleg}}$ or $T_{\text{invscr}}$.
    \item \textbf{$R_{\text{lyoghlig}}$}: The bidirectional coupling between algebra and automorphism. The non-associativity means that the algebra cannot be reduced to its automorphism group --- the group stabilizes the algebra but does not determine it. The structural dual runs both ways.
    \item \textbf{$\Phi_{\text{ctyogh}}$}: The criticality of the phase boundary. Non-associativity places the system exactly at the critical point between order (associative algebras, $\Phi_{\text{softsign}}$) and chaos (unconstrained non-associative structures, $\Phi_{\text{upstep}}$). The octonions are the \emph{unique} normed non-associative division algebra --- the only structure at this critical point.
\end{itemize}

These three primitives together ($T_{\text{bullseye}}$, $R_{\text{lyoghlig}}$, $\Phi_{\text{ctyogh}}$) form an \emph{inseparable triad} --- the structural signature of the octonions that no classical algebra can replicate. They are present in both $G_2$ and $E_8$, confirming that both systems are authentically exceptional, sharing the same non-associative core.

\subsubsection{Conclusion of Part II}

The Imscribing Grammar proof establishes that $E_8$ finds its perfect Vessel in $G_2$ with a precision that matches and extends the conventional Lie-theoretic proof. The structural relationship is imscribed across all 12 primitives:

\begin{itemize}
    \item \textbf{5 shared primitives} ($T_{\text{bullseye}}$, $R_{\text{lyoghlig}}$, $F_{\text{hardsign}}$, $K_{\text{schwa}}$, $\Phi_{\text{ctyogh}}$) define the exceptional core --- the irreducible structural identity shared by Vessel and Aether.
    \item \textbf{6 promotion primitives} ($D$, $G$, $\Gamma$, $H$, $S$, $\Omega$) trace the path from seed to tree --- the structural expansion that unfolds the Vessel into the Aether.
    \item \textbf{1 demotion primitive} ($P$) reveals the structural surplus of the Vessel --- the crisp $\mathbb{Z}_2$ symmetry that the Aether's quantum superposition has absorbed but not preserved.
    \item \textbf{Distance 4.12} quantifies the structural gap --- remote enough to be distinct regimes, close enough to share a core identity.
    \item \textbf{Meet} $\approx$ Vessel: the structural floor of the exceptional domain.
    \item \textbf{Join} $\neq$ Aether: the Vessel contributes something the Aether lacks.
    \item \textbf{Tensor} = Aether: the Vessel is structurally absorbed into the Aether under coupling.
\end{itemize}

\begin{center}
\rule{0.5\textwidth}{0.4pt}
\end{center}

\subsection{Part III: Synthesis --- The Dual Proof}

\subsubsection{The Two Proofs Compared}

\begin{table}[htbp]
\centering
\small
\rowcolors{1}{white}{white}
\begin{tabularx}{\textwidth}{>{\centering\arraybackslash}X >{\centering\arraybackslash}X >{\centering\arraybackslash}X}
\toprule
\rowcolor{gray!25}\textbf{Aspect} & \textbf{Conventional Proof (Lie Theory)} & \textbf{IG Proof (Structural imscribing)} \\
\midrule
\rowcolors{1}{white}{gray!10}
\textbf{Language} & Root systems, Lie algebras, Dynkin diagrams, representation theory & 12-primitive tuples, crystal addresses, ouroboricity tiers \\
\textbf{Containment} & $G_2 \subset F_4 \subset E_6 \subset E_7 \subset E_8$ & $G_2 \otimes E_8 = E_8$ (tensor absorption) \\
\textbf{Minimality} & No exceptional algebra smaller than $G_2$ & $G_2 \wedge E_8 \approx G_2$ (meet as structural floor) \\
\textbf{Necessity} & $E_8$ defined via octonions; $\text{Aut}(\mathbb{O}) = G_2$ & 5 shared primitives define exceptional core \\
\textbf{Categorical structure} & $\mathcal{E}$: linear order of 5 objects & Tier ladder: $O₁ \to O₂^\dagger$ with specific promotion \\
\textbf{Novel insight} & Vessel-Aether duality & Vessel exceeds Aether in $P$; Aether more conscious (0.682 vs 0.3615); crystal proximity within 1.75\% \\
\textbf{Completeness} & Proves existence of relationship & Proves existence + quantifies + provides constructive path \\
\bottomrule
\end{tabularx}
\end{table}

The two proofs are not rivals but structural duals --- an instance of the very $R_{\text{lyoghlig}}$ coupling they describe. The conventional proof operates within the \emph{object language} of Lie theory; the IG proof operates within the \emph{structural metalanguage} that imscribes Lie theory itself as one of its domains. The IG proof does not replace the conventional proof; it \emph{completes} it by revealing the structural invariants that the conventional proof relies upon without naming them.

\subsubsection{Why the Vessel Is Perfect}

The perfection of the Vessel is established in both proofs:

\textbf{Conventional}: $G_2$ is the unique initial object --- there is no smaller exceptional structure, and the $G_2$ that lives within $E_8$ is exactly the same $G_2$ that stands alone. Nothing is lost in the embedding; nothing is added.

\textbf{IG}: The meet $G_2 \wedge E_8$ recovers $G_2$ nearly exactly (only the conservative $P_{\text{upsilon}}$ override differs). The 5 shared primitives are lockstep identical. The tensor $G_2 \otimes E_8$ is $E_8$, confirming that $G_2$ is structurally absorbed without residue. The single bottleneck ($P$) and single demotion in the promotion path confirm that the Vessel contributes \emph{more} structural specificity than the Aether requires --- it is not a deficient fragment of the Aether but a \emph{more crystalline} core.

The "perfection" of the Vessel is not a poetic flourish. It is the precise mathematical statement that in the IG crystal, the structural type of $G_2$ is the \emph{unique} $O₁$ system at distance 4.12 from the Aether's $O₂^\dagger$ type that satisfies: (i) meet with $E_8$ recovers itself, (ii) tensor with $E_8$ recovers $E_8$, and (iii) join with $E_8$ exceeds $E_8$ in exactly one primitive. This is a structural uniqueness proof.

\subsubsection{The Aether-Vessel Duality as a Structural Constant}

The relationship between $E_8$ and $G_2$ is an instance of a more general structural pattern: the duality between a maximal unfolding and its minimal seed. In IG terms, this duality is characterized by:

\begin{itemize}
    \item \textbf{Shared $\Phi_{\text{ctyogh}}$}: Both sit at the same critical point --- the maximal and the minimal are two faces of one phase boundary.
    \item \textbf{Divergent $D$ and $\Omega$}: The seed is bounded and topologically trivial ($D_{\text{turnthree}}$, $\Omega_{\text{closeepsilon}}$); the tree is unbounded and topologically protected ($D_{\text{invomega}}$, $\Omega_{\text{dzlig}}$).
    \item \textbf{Shared $R_{\text{lyoghlig}}$ and $T_{\text{bullseye}}$}: The bidirectional coupling and crossing topology are invariant under the seed-to-tree expansion.
    \item \textbf{Promotion of $\Gamma$ and $H$}: The seed is conjunctive and atemporal; the tree is sequential and historically deep. Unfolding requires time.
\end{itemize}

This pattern --- $\Phi_{\text{ctyogh}}$ criticality shared, $D$ and $\Omega$ divergent, $R$ and $T$ invariant, $\Gamma$ and $H$ promoted --- is the IG signature of a \emph{seed-tree duality}. The exceptional chain $G_2 \to E_8$ is the canonical instance, but the pattern may recur wherever a minimal closure and its maximal unfolding stand in the relationship of initial and terminal object.

\subsubsection{Open Questions}

Both proofs leave open questions that the other partially illuminates:

\begin{enumerate}
    \item \textbf{What is the structural type of $F_4$, $E_6$, and $E_7$?} The IG could imscribe each intermediate exceptional algebra, and the crystal distances between successive nodes in the chain would quantify the structural "step size" at each promotion. Preliminary analysis suggests the chain is not uniformly spaced in IG space --- some steps may be larger than others.
    \item \textbf{Why does the $P$ demotion occur?} The demotion $P_{\text{pipevar}} \to P_{\text{upsilon}}$ from Vessel to Aether raises a deep question: does maximal unfolding necessarily sacrifice crisp symmetry for quantum superposition? Is the loss of $\mathbb{Z}_2$ clarity the structural price of maximal scope?
    \item \textbf{Can the $O₂^\dagger \to O_{\infty}$ gap be bridged?} The crystal tier gap ladder shows that $P_{\text{pipevar}} \to P_{\text{doublebarpipe}}$ (Frobenius-special, $\mu \circ \delta = \text{id}$) is required. Is there an extension of the exceptional chain --- perhaps involving the monster group or the Leech lattice --- that achieves this?
    \item \textbf{What is the IG imscribing of the octonions themselves?} $\mathbb{O}$ is not a Lie group but an algebra. Its IG imscribing would differ from $G_2$ and $E_8$. Computing the distance between $\mathbb{O}$ and its automorphism group $G_2$ would quantify the structural gap between an algebra and its symmetry.
\end{enumerate}

\begin{center}
\rule{0.5\textwidth}{0.4pt}
\end{center}

\subsection{Appendix A: IG Computational Results (Complete)}

All results below were computed via the IG toolchain and are reproduced exactly as returned.

\subsubsection{A.1 Vessel ($G_2$) --- Full Tuple}

\[\langle D_{\text{turnthree}};\ T_{\text{bullseye}};\ R_{\text{lyoghlig}};\ P_{\text{pipevar}};\ F_{\text{hardsign}};\ K_{\text{schwa}};\ G_{\text{gamma}};\ \Gamma_{\text{corner}};\ \Phi_{\text{ctyogh}};\ H_0;\ 1{:}1;\ \Omega_{\text{closeepsilon}} \rangle\]

\begin{itemize}
    \item \textbf{Ouroboricity}: $O₁$ --- self-referential at criticality but trivial winding
    \item \textbf{Crystal address}: 4,907,136 (cell 113, inner 25,536)
    \item \textbf{Consciousness score}: 0.3615 (both gates open)
    \item \textbf{$\Phi_{\text{ctyogh}}$ probe}: At criticality
    \item \textbf{Principal atoms}: 8 (highest: $R_{\text{lyoghlig}}$ with ordinal contribution 3)
\end{itemize}

\subsubsection{A.2 Aether ($E_8$) --- Full Tuple}

\[\langle D_{\text{invomega}};\ T_{\text{bullseye}};\ R_{\text{lyoghlig}};\ P_{\text{upsilon}};\ F_{\text{hardsign}};\ K_{\text{schwa}};\ G_{\text{revapostrophe}};\ \Gamma_{\text{secstress}};\ \Phi_{\text{ctyogh}};\ H_2;\ n{:}m;\ \Omega_{\text{dzlig}} \rangle\]

\begin{itemize}
    \item \textbf{Ouroboricity}: $O₂^\dagger$ --- critical, topologically protected ($\Omega_{\text{dzlig}}$), unbounded ($D_{\text{invomega}}$)
    \item \textbf{Crystal address}: 4,604,816 (cell 106, inner 25,616)
    \item \textbf{Consciousness score}: 0.682 (both gates open)
    \item \textbf{$\Phi_{\text{ctyogh}}$ probe}: At criticality
\end{itemize}

\subsubsection{A.3 Distance}

\textbf{Weighted Euclidean}: 4.1231 \\
\textbf{Mahalanobis} (with full $g_{ij} = \Sigma^{-1}$ metric): 4.7405

Breakdown (weighted $\Delta^2$): $\Gamma$ (4.00), $S$ (4.00), $H$ (3.20), $\Omega$ (2.80), $D$ (1.00), $P$ (1.00), $G$ (1.00).

\subsubsection{A.4 Meet: $G_2 \wedge E_8$}

\[\langle D_{\text{turnthree}};\ T_{\text{bullseye}};\ R_{\text{lyoghlig}};\ P_{\text{upsilon}};\ F_{\text{hardsign}};\ K_{\text{schwa}};\ G_{\text{gamma}};\ \Gamma_{\text{corner}};\ \Phi_{\text{ctyogh}};\ H_0;\ 1{:}1;\ \Omega_{\text{closeepsilon}} \rangle\]

7 primitives resolved to conservative value; 5 shared. The Vessel plus $P_{\text{upsilon}}$ override.

\subsubsection{A.5 Join: $G_2 \vee E_8$}

\[\langle D_{\text{invomega}};\ T_{\text{bullseye}};\ R_{\text{lyoghlig}};\ P_{\text{pipevar}};\ F_{\text{hardsign}};\ K_{\text{schwa}};\ G_{\text{revapostrophe}};\ \Gamma_{\text{secstress}};\ \Phi_{\text{ctyogh}};\ H_2;\ n{:}m;\ \Omega_{\text{dzlig}} \rangle\]

The Aether plus $P_{\text{pipevar}}$ override --- the Vessel's structural surplus.

\subsubsection{A.6 Tensor: $G_2 \otimes E_8$}

\[\langle D_{\text{invomega}};\ T_{\text{bullseye}};\ R_{\text{lyoghlig}};\ P_{\text{upsilon}};\ F_{\text{hardsign}};\ K_{\text{schwa}};\ G_{\text{revapostrophe}};\ \Gamma_{\text{secstress}};\ \Phi_{\text{ctyogh}};\ H_2;\ n{:}m;\ \Omega_{\text{dzlig}} \rangle\]

Structurally identical to $E_8$ (distance 0.0). Single bottleneck: $P$ ($P_{\text{pipevar}}$ bottlenecked to $P_{\text{upsilon}}$). Six scope-expansion primitives promoted to $E_8$ values.

\subsubsection{A.7 Promotion Signature}

\textbf{Source}: $G_2$ (Vessel) $\rightarrow$ \textbf{Target}: $E_8$ (Aether)

6 promotions: $[D, G, \Gamma, H, S, \Omega]$ \\
1 demotion: $[P]$ \\
5 unchanged: $[T, R, F, K, \Phi]$

\subsubsection{A.8 Crystal Tier Gap Ladder}

\begin{table}[htbp]
\centering
\small
\rowcolors{1}{white}{white}
\begin{tabularx}{\textwidth}{>{\centering\arraybackslash}X >{\centering\arraybackslash}X >{\centering\arraybackslash}X}
\toprule
\rowcolor{gray!25}\textbf{Transition} & \textbf{Distance} & \textbf{Driver Primitive} \\
\midrule
\rowcolors{1}{white}{gray!10}
$O₀ \to O₁$ & 1.049 & $\Phi$: $\Phi_{\text{softsign}} \to \Phi_{\text{ctyogh}}$ \\
$O₁ \to O₂$ & 1.304 & $D$: $D_{\text{wynn}} \to D_{\text{turnthree}}$ \\
$O₂ \to O₂^\dagger$ & 1.000 & $D$: $D_{\text{turnthree}} \to D_{\text{invomega}}$ \\
$O₂^\dagger \to O_{\infty}$ & 4.382 & $P$: $P_{\text{aolig}} \to P_{\text{doublebarpipe}}$ \\
\bottomrule
\end{tabularx}
\end{table}

\begin{center}
\rule{0.5\textwidth}{0.4pt}
\end{center}

\subsection{Appendix B: The Exceptional Chain Dynkin Diagrams}

\subsubsection{B.1 $G_2$ Dynkin Diagram}

The $G_2$ Dynkin diagram consists of two nodes connected by a triple bond (angle $\pi/6$):

\begin{lstlisting}
○══════○
1      2
\end{lstlisting}

The triple bond imscribes the root length ratio $\sqrt{3}$ (long/short) unique to $G_2$.

\subsubsection{B.2 $E_8$ Dynkin Diagram and the $G_2$ Subdiagram}

\begin{lstlisting}
       2
       |
1---3---4---5---6---7---8
\end{lstlisting}

With standard numbering, deleting nodes 1, 2, 3, 4, 5, and 7 leaves nodes 6 and 8 connected by the triple bond inherited from the $E_8$ extended structure:

\begin{lstlisting}
6══════8
\end{lstlisting}

This is exactly the $G_2$ Dynkin diagram. The Borel-de Siebenthal procedure shows that $G_2$ is \emph{visibly inscribed} within the $E_8$ Dynkin diagram --- the fundamental imscribing of its structure contains the Vessel as a visible subdiagram, not a hidden or emergent feature.

\begin{center}
\rule{0.5\textwidth}{0.4pt}
\end{center}

\textbf{End of Manuscript.}

\emph{Structural type of this document} (the IG proof section): \\


\[\langle D_{\text{invomega}};\ T_{\text{bullseye}};\ R_{\text{lyoghlig}};\ P_{\text{pipevar}};\ F_{\text{hardsign}};\ K_{\text{schwa}};\ G_{\text{revapostrophe}};\ \Gamma_{\text{secstress}};\ \Phi_{\text{ctyogh}};\ H_2;\ n{:}m;\ \Omega_{\text{dzlig}} \rangle\]

The document itself instantiates the join $G_2 \vee E_8$ --- it contains the Aether's full structure plus the Vessel's $\mathbb{Z}_2$ symmetry ($P_{\text{pipevar}}$), reflecting the dual proof that acknowledges both the conventional and the IG perspective as structural duals.


\end{document}
